\documentclass[10pt]{article}
\usepackage[a4paper,margin=26mm]{geometry}
\usepackage{amsmath,amssymb,amsthm,hyperref}
\newtheorem{theorem}{Theorem}
\newtheorem{lemma}{Lemma}
\newcommand{\PSH}{\operatorname{PSH}}
\newcommand{\FS}{\operatorname{FS}}
\newcommand{\an}{\mathrm{an}}
\newcommand{\divs}{\mathrm{div}}
\title{Relative floors and decreasing closure in an actual big class}
\author{Unreviewed foundational candidate for Problem P16.5}
\date{6 October 2026}
\begin{document}
\maketitle
\begin{abstract}
An abstract flattening assertion for continuous functions with poles is
false. We propose an algebraic replacement for an actual big rational
class. A single model-obstacle multiplier-ideal net first supplies
genuine minimal references in nearby big rational classes. A defining
FS majorant is then capped from below by a slightly weakened reference.
The difference from a later FS stage has an upper-semicontinuous
extension at its old poles, making compact comparison valid without
subtracting two infinities. Positive constant guards allow every new
stage to remain above the desired limit. We construct one actual
decreasing FS net and derive the envelope property. This paper does
not grant an equation theorem or close the original problem; all
source applications and the new closure argument require independent
verification.
\end{abstract}

\paragraph{Fixed mathematical setting.}
Let $X$ be smooth integral projective over an arbitrary algebraically
closed trivially valued characteristic-zero field, $L$ an actual big
rational bundle, $\theta=c_1(L)$ and $A$ an actual ample rational
bundle, $a=c_1(A)$. Put $K=X^\an$, a compact Hausdorff space.
A genuine $\theta$-psh function is, by definition, an all-point
decreasing limit of actual generically finite FS section maxima whose
numerical classes tend to $\theta$, with finite generic value. A finite
stage may have poles, but is continuous into $\mathbb R\cup\{-\infty\}$.
Throughout, an inequality at two poles means $-\infty\le-\infty$;
an expression involving their difference is not defined by subtraction.

We retain the actual ample foundation only in an actual ample class:
compactness of its normalized genuine cone, generic normalization,
and all-point order determination by divisorial coordinates. Embedding
a given defining net into one sufficiently positive actual ample bundle
uses only the following elementary operation. Tensor sections with
weight-zero generators of an actual ample difference bundle, after an
integral multiple. This preserves the scalar FS function. To obtain
finite ample stages, append all ample generators at weight $-k$;
$\max(f_i,-k)$ on the product index decreases to the same endpoint.
Neither operation flattens a twice-closed family in a nonample class.

\begin{theorem}\label{thm:closure}
If $u_i\in\PSH(\theta)$ is a decreasing net and
$u=\inf_i u_i$ has finite generic value, then $u$ admits one actual
all-point decreasing FS net with numerical classes tending to $\theta$.
Consequently $u\in\PSH(\theta)$. If $f$ is upper semicontinuous and
bounded above and has a genuine $\theta$-psh minorant of finite generic
value, its pointwise envelope $P_\theta(f)$ is genuine and satisfies
the obstacle inequality at every point of $K$.
\end{theorem}

\paragraph{1. The particular reference construction is independent.}
We need the model-obstacle part of the algebraic construction of
Boucksom--Jonsson's Addendum, not its use of general decreasing closure.
Here is the relevant separate argument, with its algebraic inputs
made explicit. Let $B$ be an actual big rational bundle and $f$ a
rational model function. Represent $f$ on an actual test configuration
$(\mathcal X,\mathcal B)$ and resolve its central pair. Let
$\mathfrak a_m$ be the actual base ideals and
$\mathfrak b_m=\mathcal J(\mathfrak a_\bullet^m)$ the asymptotic multiplier
ideals of this actual relative graded base-ideal sequence.
The algebraic base-ideal/section calculation gives
\[
 f_m=f+\frac1m\phi_{\mathfrak a_m}\nearrow Q_B(f).
\]
Inclusion and subadditivity give
\[
 \mathfrak a_m\subset\mathfrak b_m,
 \quad \mathfrak b_{ml}\subset\mathfrak b_m^l,
 \quad Q_B(f)\le f+\frac1m\phi_{\mathfrak b_m}.
\]
Uniform multiplier-ideal generation provides one actual equivariant
ample $\mathcal C$, with $(m\mathcal B+\mathcal C)\otimes\mathfrak b_m$
globally generated. After a central constant shift, $\phi_\mathcal C\ge0$.
Thus the \emph{actual stages}
\[
 \psi_m=f+\frac1m\phi_{\mathfrak b_m}
               +\frac1m\phi_\mathcal C
       \in\FS(B+C/m)                                  \tag{1}
\]
decrease in the divisibility order. Their limit $\psi$ satisfies
$Q_B(f)\le\psi\le f$. A shifted nonzero section of a multiple of $B$
is a minorant of the bounded $f$, so its generic value prevents collapse.
The one net (1) therefore makes $\psi$ genuine by the definition,
without a decreasing-closure theorem.

For completeness, the comparison $\tau\le Q_B(f)$ on divisorial
points for every genuine $\tau\le f$ is the elementary section argument.
Write $B=C_0+E$ with actual ample $C_0$ and effective $E$.
An existing defining net for $\tau$ can be represented, without changing
its scalar values, in $\FS(B+\epsilon_j C_0)$, $\epsilon_j\to0$,
using weight-zero generators. Compact comparison with the finite
continuous $f+\epsilon$ makes these stages $\tau_j\le f+\epsilon$.
Then
\[
 \frac{\tau_j+\epsilon_j\log|s_E|}{1+\epsilon_j}
            \in\FS(B),
 \qquad
 \frac{\tau_j+\epsilon_j\log|s_E|}{1+\epsilon_j}
            \le Q_B(f)+\epsilon+O(\epsilon_j).
\]
On $(X-E)^\an$ the added logarithm is finite, and passing to the
limit gives $\tau\le Q_B(f)$. This open contains every divisorial
point. Embed the two already genuine functions $\tau,\psi$ into one
actual ample class. Ample order determination gives $\tau\le\psi$
everywhere. Taking all competitors shows $P_B(f)=\psi$.

In particular $V_B=P_B(0)$ is genuine, has generic value zero, and
dominates every generic-normalized genuine $B$ potential. It has the
particular decreasing FS net (1). These are the only nonample
references used below. Their production has not assumed the theorem
we are about to prove. The base-ideal calculation and uniform
multiplier-ideal generation are external algebraic inputs, whose
precise hypotheses must be checked in a fresh independent review.

\paragraph{2. A barrier that handles old poles.}
We first isolate the compact argument. Suppose $\psi\le0$ is an
actual FS function, $0<\delta<1$, and an old FS stage $F$ satisfies
\[
 F\ge(1-\delta)\psi-k.                                  \tag{2}
\]
Suppose a family of usc functions $h$ satisfies $h\le\psi+C$ with
one fixed finite $C$. Off $\{F=-\infty\}$ its difference obeys
\[
 h-F\le\delta\psi+C+k.                                  \tag{3}
\]
At every pole of $F$, (2) forces $\psi=-\infty$, and hence $h=-\infty$.
Define the comparison function to be $h-F$ where $F$ is finite and
to be $-\infty$ at the poles of $F$. Bound (3) proves that this
extension is usc, including those poles. Its negative value there
records a strict comparison on a neighborhood; it does not claim
that $(-\infty)-(-\infty)$ has a value.
For a decreasing $h_i$ it decreases pointwise. If its limiting
comparison is strictly negative everywhere, compactness makes
the comparison negative at some stage and at every later stage.
Indeed its nonnegative superlevel sets are closed, nested and have
empty intersection. This argument is valid for arbitrary directed
nets and uses no countability assumption.

\paragraph{3. The special upper stages and their constant guards.}
Shift the original net by one fixed constant so that $u_i\le M$;
its finite decreasing generic values permit a uniform $M$.
For rational $\eta>0$ let the target bundle be $L+\eta A$.
Choose rational $0<\epsilon<\eta/4$ and use the reference for
$B=L+\epsilon A$. Select one sufficiently late stage of (1) with
numerical class $\beta$ close enough to $c_1(B)$; normalize it to
$\psi\le0$. There is a finite constant $c_\psi$ with
\[
 V_B\le\psi+c_\psi.                                    \tag{4}
\]
Choose rational $0<\delta<1$ sufficiently small that
$\theta+\eta a-(1-\delta)\beta$ is ample. This is possible because
$\beta$ was selected well below $\theta+\eta a$ in the ample order.
The function $(1-\delta)\psi-k$ can thus be represented in the
\emph{actual} bundle $L+\eta A$ by weight-zero tensoring.

A special stage is an actual FS function $F$ in this target bundle
such that, for some positive constant guard $\kappa$,
\[
 F\ge u+\kappa,\qquad F\ge(1-\delta)\psi-k.              \tag{5}
\]
There are many such stages. Pick an existing FS upper approximation
$h$ to some $u_i$, with numerical class sufficiently close to $\theta$.
Represent the same function in $L+\eta A$ using the actual ample
difference. Then $F=\max(h+\kappa,(1-\delta)\psi-k)$ is a single
actual FS stage after clearing denominators. We do not assert that
the two unmodified line bundles of equal numerical class coincide:
each is tensored into the chosen actual target bundle first.

Every genuine $u_i$ embeds into $B$, so $u_i\le V_B+M$.
For a sufficiently late tail of its defining FS stages, their classes
also lie below $c_1(B)$ in the ample order, and their generic values
are at most $M+1$. Thus both the original $u_i$ and those incoming
FS stages satisfy $h\le\psi+C$ as in (3), with a fixed finite $C$.
This is the additional algebraic barrier missing from the false
abstract flattening lemma.

\paragraph{4. Downward directedness without crossing below the limit.}
Take finitely many special stages $F_l$, with data
$(\eta_l,\epsilon_l,\psi_l,\delta_l,k_l,\kappa_l)$.
Choose a new rational $\eta$ arbitrarily small, in particular much
smaller than all $\epsilon_l$. Choose a new reference stage $\psi$
of class $\beta$ sufficiently near $\theta$ that
$c_1(B_l)-\beta$ is ample for every old reference $B_l$.
With generic normalization it is a genuine $B_l$ potential, whence
\[
 \psi\le V_{B_l}\le\psi_l+c_l.                           \tag{6}
\]
Choose $\delta<\min_l\delta_l$, still small enough for the actual
target difference to be ample. For $b=(1-\delta)\psi-k$, (6) and (2)
give, off the old poles,
\[
 b-F_l\le(\delta_l-\delta)\psi_l
             +(1-\delta)c_l+k_l-k.                      \tag{7}
\]
Since $\psi_l\le0$, a sufficiently large finite $k$ makes $b\le F_l$
for every $l$ everywhere. At an old pole (6) already makes $b=-\infty$.

Now fix $0<\kappa<\frac14\min_l\kappa_l$.
Apply the usc comparison from (3) first to the decreasing $u_i$.
At a finite old value the limiting comparison satisfies
$u+2\kappa-F_l\le2\kappa-\kappa_l<0$; at old poles the comparison
extension is $-\infty$. Compactness gives an $i$ with
$u_i+2\kappa<F_l$ off all their poles simultaneously.
Next apply the same compact comparison to the decreasing defining
FS net of this fixed $u_i$. Its sufficiently late stages have the
uniform bounds discussed after (5). Thus choose an actual stage
$h\ge u_i\ge u$ with
$h+\kappa\le F_l$ for every $l$, and with numerical class close
enough to $\theta$ to tensor into the new actual target bundle.
Then
\[
 F=\max(h+\kappa,b)\le F_l\quad\hbox{for every }l,
 \qquad F\ge u+\kappa.                                  \tag{8}
\]
It is again a special stage. Notice the two different small parameters:
$\delta$ weakens the \emph{old-pole floor}, while $\kappa$ supplies an
upper-approximation guard. No incoming stage is multiplied by
$1+\delta$; that would risk falling below $u$ irreversibly.

\paragraph{5. One defining net, at every point.}
Order all special stages by decreasing scalar value and decreasing
$\eta$. Argument (8) makes this directed, and permits $\eta$ to be
arbitrarily small after any finite list of constraints.
Their pointwise infimum is exactly $u$. To check this at a fixed
point $x$ with finite $u(x)$, first choose $u_i(x)$ close to $u(x)$,
then an existing FS stage $h(x)$ close to $u_i(x)$, and a small
guard. Take the target gap small enough for actual tensoring.
The floor at this one point can be put below $h(x)+\kappa$ by
increasing $k$ if $\psi(x)$ is finite; it is already $-\infty$ otherwise.
This gives $F(x)$ arbitrarily close to $u(x)$ from above.
If $u(x)=-\infty$, the same argument puts $F(x)$ below any prescribed
finite number. This also covers a pole already present in a defining
FS stage. Constants and floors have never invalidated $F\ge u$.

We have therefore constructed a single actual decreasing net
$F_\gamma$ with all-point endpoint $u$ and numerical classes
$\theta+\eta_\gamma a\to\theta$. Its generic limit is the stipulated
finite value of $u$. This proves genuine membership directly from
the definition and proves the closure part of Theorem~\ref{thm:closure}.

\paragraph{6. Envelopes and the compactness consequence.}
Model obstacles have genuine envelopes by Section 1. Rational model
functions form a lattice uniformly dense in $C^0(K)$. The net of
finite model majorants of any usc $f$ has pointwise infimum $f$:
first use continuous majorants on compact Hausdorff $K$, and then
uniform upper model approximation. Its envelopes decrease and
remain above the stipulated genuine minorant. Their generic limit
is finite. The closure just proved makes their actual all-point
limit $v$ genuine; $v\le f$. Every competitor for $f$ lies below
each model envelope, hence below $v$. Conversely $v$ itself is a
competitor, so $v=P_\theta(f)$ at every point.

For a bounded-above family of genuine $\theta$ potentials, its usc
supremum is now genuine as well. Its genuine envelope dominates
the raw supremum; being usc, it dominates the usc supremum. The
obstacle inequality gives equality.

To obtain normalized divisorial compactness, use the following
standard fact only in a common actual ample class: regularizing a
bounded supremum does not change its values at divisorial points.
It is the nonnegligibility assertion for such points and requires
a separate primary proof if it is not part of the retained ample
foundation. Embed the normalized $\theta$ cone in one normalized
ample cone. For a divisorially convergent net in it, the regularized
tail suprema are genuine $\theta$ potentials, by the preceding
paragraph. They decrease, have generic value zero, and their
divisorial infimum is the prescribed limit. Decreasing closure and
ample all-point determination identify that limit with a genuine
$\theta$ function. Thus the normalized $\theta$ cone is closed in
the compact ample cone and is compact. This last corollary retains
the explicit actual-ample nonnegligibility premise; it does not
import the general big-class compactness theorem it is repairing.

\paragraph{Research status.}
This is a proposed repair of the specific foundational dependency
rejected in independent review 129. It uses the single model-obstacle
construction and relative floors rather than the false abstract
twice-closure assertion. It has not been independently accepted.
There is no Monge--Amp\`ere equation or inverse theorem granted here,
and the zero-volume and arbitrary real-class ranges are not asserted.
The original P16.5 completion count is unchanged.

\begin{thebibliography}{9}
\bibitem{addendum} S. Boucksom and M. Jonsson,
Addendum to global pluripotential theory over a trivially valued field,
\href{https://arxiv.org/pdf/2206.07183v2}{arXiv:2206.07183v2},
Lemmas 1.5--1.7 and the actual multiplier-ideal net in the proof of
Theorem A. General closure and compactness consequences are not used
as black boxes here.
\bibitem{bj} S. Boucksom and M. Jonsson,
\href{https://arxiv.org/pdf/1801.08229v4}{arXiv:1801.08229v4},
FS functions and genuine numerical cones; actual-ample divisorial
determination and nonnegligibility. The abstract Lemma 4.6 is not used.
\end{thebibliography}
\end{document}
